\input{praeambel}
\usepackage{longtable}
\usepackage{array}
\usepackage{booktabs} % for nicer lines

\usepackage{titling}
\setlength{\droptitle}{-10em}

\newenvironment{timeline}{%
    \begin{longtable}{@{} p{1.5cm} p{9cm} p{2.5cm} @{}}
    \multicolumn{3}{l}{\textbf{Timeline}} \\
    \toprule
      \textbf{Datum} & \textbf{Aktivität} & \textbf{Aufwand} \\
    \midrule
    \endfirsthead
    \multicolumn{3}{l}{\textbf{Timeline (Fortsetzung)}} \\
    \toprule
    \textbf{Datum} & \textbf{Aktivität} & \textbf{Aufwand} \\
    \midrule
    \endhead
    \midrule
    \multicolumn{3}{r}{{Weiter nächste Seite}} \\
    \endfoot
    \bottomrule
    \endlastfoot
}{%
    \end{longtable}
}

% Center the title page
\newgeometry{hmarginratio=1:1}

\title{Diagrams as a Service \\ \small Forschungstagebuch}
\date{31. März 2026}
\author{Henri Heyden \\%
 \small stu240825 \\%
 \small Christian-Albrechts-Universität zu Kiel%
}

\begin{document}
    \maketitle
    
    \noindent
    Für mein Forschungsprojekt im Wintersemester 2025/26, habe ich wie nach folgender Timeline gearbeitet.
    Es ist empfehlenswert für Kontext vorerst den Bericht über das Forschungsprojekt zu lesen.

    \begin{timeline}
        30.10.25 & Beginn des Forschungsprojekts; Onboarding in die KIELER Infrastruktur. & 6 Stunden \\

        6.11.25 & Aufsetzen und Konfiguration von Git repositories und VS Code workspaces um am Projekt arbeiten zu können; Aufsetzen von KGT-Server ohne Microlayout sowie KGT Beispielen; Fixes um simple Beispiele ohne Abstürze laufen zu lassen; nötige Nullchecks; Heuristik um Textgrößen abschätzen zu können; Anpassung von Client Klassen und Typprüfungen; Start Implementierung Microlayout Abschätzung. & 6 Stunden \\

        14.11.25 & Einführende Präsentation. & 6 Stunden \\

        19.11.25 & Aufsetzen von Server Quellcode in Eclipse um Porten leichter zu machen; Implementierung von diversen internen Funktionen für Renderings die flächenplatziert sind. & 5 Stunden \\

        20.11.25 & Umformatieren und Säubern von Code; Korrigieren einiger Typprüfungen; Implementierung von anderen Funktionen für flächenplatzierte Renderings und relevante Typprüfungen; Fixes; Warnungen für ungetestete Methoden; Implementierung von wichtigen Funktionen für punktplatzierte Renderings. & 5 Stunden \\

        26.11.25 & Navigation im Server Code um Implementierung der Textgrößen nachvollziehen zu können; Implementierung weiterer Funktionen für punktplatzierte Renderings; Implementierung der Abschätzung von Größen von Bildern (\lstinline|KImage|); Start der Implementierung der Einschätzung von rasterplatzierten Renderings und allgemeinen wichtigen Funktionen; Fixes; weitere Einschätzung der Implementierung. & 7 Stunden \\

        11.12.25 & Start der Implementierung von Microlayout Berechnung für ersten Durchstich; Fixes; Weitere Arbeit an rasterplatzierten Renderings. & 4 Stunden \\

        7.1.26 & Planung der weiteren Implementierung; Implementierung von vielen wichtigen Funktionen um die Größe von Renderings nach dem Makrolayout zu berechnen. & 4 Stunden \\

        15.1.26 & Weitere Planung und Priorisierung der weiteren Implementierung; Berücksichtigung von punktplatzierten Renderings im Microlayout-Berechnungs Schritt; Fixes und Säubern des Codes; Implementierung der Microlayout Berechnung von Renderings innerhalb Renderings (\lstinline|KContainerRenderings|). & 5 Stunden \\

        20.1.26 & Säubern des Codes für die Microlayout Berechnung; Implementierung der Navigation im \lstinline|SKGraph| für Knoten und Labels; Fixes & 5 Stunden \\

        22.1.26 & Viel Säubern und viele Fixes; Teilweise Implementierung der Microlayout Berechnung Renderings welche mithilfe von \lstinline|KDecoratorPlacementData| plaziert werden. & 4 Stunden \\
        
        24.1.26 & Weiteres Säubern und wichtige Fixes; Planung wie die minimale Größe von \lstinline|KNodes| während der Microlayout Schätzung herausgefunden wird; nachdenken wie zwischen verschiedenen Kind-Typen einer \lstinline|KNode| in der Microlayout Berechnung unterschieden werden muss; Versuch das generierte Padding in einem Graphen zu entfernen & 6 Stunden \\

        25.1.26 & Implementierung der \lstinline|KDecoratorPlacementData| Klasse und einer dazugehörigen Typprüfung. & 1 Stunde \\

        26.1.26 & Vorbereitung und Präsentation der Zwischenergebnisse. & 5 Stunden \\

        29.1.26 & Fixes. & 2 Stunden \\

        4.2.26 & Typprüfung für \lstinline|KRenderingLibrary| Klasse; Implementierung einer Suchfunktion, die in einem Graphen die jeweilige \lstinline|KRenderingLibrary| zurückgibt. & 2 Stunden. \\

        6.2.26 & Anfang Implementierung einer Test-Suite um Client Microlayout Funktionen anhand von Beispielen zu testen und mit Server Microlayout zu vergleichen; erfolgreiches Entfernen von Padding; Anfang einer Implementierung der \lstinline|KRenderingRef| Resolution; Fixes. & 5 Stunden \\

        23.2.26 & Arbeit am Forschungstagebuch. & 6 Stunden \\

        24.2.26 & Pull-request; Arbeit am Forschungsbericht & 6 Stunden \\
    \end{timeline}
    
    \noindent
    Somit ergibt sich insgesamt ein Zeitaufwand von 90 Stunden.
    Es sei angemerkt, dass sehr viel Debugging und Arbeit anhand Beispielen nötig war, weswegen einige Zeiten unproportional zu den Anmerkungen vorkommen können.

    Die wesentliche Arbeit in dem Forschungsprojekt lässt sich auch anhand meiner Änderungen in dem folgenden Repository vermerken: Klighd-Vscode (\url{https://github.com/kieler/klighd-vscode/tree/nre/henri-fork}).
    Die einführende Präsentation sowie die Zwischenpräsentation sind zu finden unter \url{https://docs.google.com/presentation/d/1q_nDqAl1FoWcPzSdqEaZ6JHt0WQSJ5WqP9HgScHuBlg/}.    
\end{document}

